\section{Aufgaben}


\Title{Multiple Choice}

Welche der folgenden Aussagen ist korrekt für $f : \R^n \mapsto \R^m, g : \R^m
\mapsto \R$ und $x \in \R^n$?
\begin{itemize}
    \item[$\square$] $f$ ist stetig in $x$, falls es eine Folge $(x_k)_{k \in \N} \subseteq
    \R^n$ gibt, so dass $f(x_k) \to f(x)$ für $k \to \infty$.
    \item[$\square$] $f \circ g$ und $g \circ f$ sind immer definiert
    \item[$\square$] falls $g \circ f$ und $f$ stetig sind, dann ist auch $g$ stetig
\end{itemize}

\hrulefill

Welche der folgenden Funktionen sind stetig?
\begin{itemize}
    \item[\rlap{\raisebox{0.3ex}{\hspace{0.4ex}\tiny \ding{52}}}$\square$] $f(x,y,z) = x^3 + 
    y^2 + 3z$
    \item[$\square$] $f(x,y) = \inf \{x^k + y^k | k \in \Z \}$, für $(x,y) \in \R^2_{\geq 0}$
    \item[\rlap{\raisebox{0.3ex}{\hspace{0.4ex}\tiny \ding{52}}}$\square$] $f(x,y) = \inf 
    \{x^k + y^k | k \in \Z \}$, für $0 \leq x,y < 1$
    \item[\rlap{\raisebox{0.3ex}{\hspace{0.4ex}\tiny \ding{52}}}$\square$] $f(x,y) = \int_x^y
    \cos(t)dt$, für $x < y$
\end{itemize}

\hrulefill

Welche der folgenden Aussagen ist wahr?
\begin{itemize}
    \item[$\square$] sei pr:$\R^2 \mapsto \R$, pr$(x,y) := x$ und $A \subseteq \R^2$
    abgeschlossen, dann ist auch pr$(A)$ abgeschlossen
    \item[\rlap{\raisebox{0.3ex}{\hspace{0.4ex}\tiny \ding{52}}}$\square$] Falls $A \subseteq
    \R^n$ geschlossen ist, dann ist $A^c$ offen
    \item[$\square$] Die Menge $O \subseteq R^n$ ist offen genau dann wenn $O$ nicht 
    abgeschlossen ist
    \item[\rlap{\raisebox{0.3ex}{\hspace{0.4ex}\tiny \ding{52}}}$\square$] Sei $f:[0,1]^n 
    \mapsto \R$ stetig, dann ist $f$ beschränkt
    \item[$\square$] Sei $f:[-1,1]^n \backslash \{0\} \mapsto \R$ stetig, dann ist $f$
    beschränkt
\end{itemize}

\hrulefill

Sei $X \subseteq \R^n$ offen und $f,g: X \mapsto \R^m$. Welche Aussagen sind korrekt?
\begin{itemize}
    \item[\rlap{\raisebox{0.3ex}{\hspace{0.4ex}\tiny \ding{52}}}$\square$] falls $f \in C^2$
    auf $X$ ist, dann ist $f$ auch $C^1$ auf $X$
    \item[$\square$] falls $f$ und $g$ jeweils $C^3$ und $C^2$ auf $X$ sind, dann ist
    $f \cdot g$ $C^3$ auf $X$
    \item[\rlap{\raisebox{0.3ex}{\hspace{0.4ex}\tiny \ding{52}}}$\square$] Sei $f = (f_1, 
    ..., f_m)$, wobei $f_i$ Polynome sind und $g$ ist $C^k$, dann ist $f/g$ $C^k$ auf $X$
    \item[\rlap{\raisebox{0.3ex}{\hspace{0.4ex}\tiny \ding{52}}}$\square$] Sei $f$ $C^k$ mit 
    $k \in \N$, dann gilt im Allgemeinen $\partial_{xy} f = \partial_{yx} f$.
\end{itemize}

\hrulefill

Sei $f: \R^n \mapsto \R$ eine $C^k$ Abbildung mit $k \geq 2$. Welche Aussagen sind korrekt?
\begin{itemize}
    \item[$\square$] der Gradient $\nabla f(x)$ ist eine $n \cross n$ Matrix
    \item[\rlap{\raisebox{0.3ex}{\hspace{0.4ex}\tiny \ding{52}}}$\square$] $H_f(x)$ ist eine
    quadratische Matrix
    \item[\rlap{\raisebox{0.3ex}{\hspace{0.4ex}\tiny \ding{52}}}$\square$] $H_f(x)$ ist 
    symmetrisch
    \item[$\square$] $H_f(x)$ ist invertierbar
\end{itemize}

\hrulefill

Sei $f: X \mapsto \R, g: K \mapsto \R$ differenzierbar, $X \subseteq \R^n$ offen und $K
\subseteq \R$ kompakt. Welche Aussagen sind korrekt?
\begin{itemize}
    \item[$\square$] $f$ besitzt eine Extremalstelle
    \item[$\square$] $g$ besitzt keine Extremalstelle
    \item[\rlap{\raisebox{0.3ex}{\hspace{0.4ex}\tiny \ding{52}}}$\square$] $g$ besitzt
    mindestens zwei Extremalstellen
    \item[$\square$] die Stelle $x_0$ ist eine lokale Minimalstelle genau dann wenn für jedes
    $\delta > 0$ mit $C(x_0, \delta) \subseteq X$ gilt: $f(y) \leq f(x_0), \forall y \in 
    C(x_0, \delta)$
    \item[$\square$] falls $df(x_0) = 0$ gilt, ist $f(x_0)$ entweder ein lokales Minimum oder
    Maximum
    \item[$\square$] die Menge der kritischen Punkte von $f$ ist immer endlich
\end{itemize}

\hrulefill

Sei $f: X \mapsto \R^n$ ein Vektorfeld und $\gamma: [a,b] \mapsto \R^n$ eine parametrisierte 
Kurve. Welche Aussagen sind korrekt?
\begin{itemize}
    \item[\rlap{\raisebox{0.3ex}{\hspace{0.4ex}\tiny \ding{52}}}$\square$] das Wegintegral
    $\int_\gamma f(s)ds$ ist eine reelle Zahl
    \item[\rlap{\raisebox{0.3ex}{\hspace{0.4ex}\tiny \ding{52}}}$\square$] falls $\gamma(t)
    \equiv v$ für alle $t \in [a,b]$, dann gilt $\int_\gamma f(s)ds = 0$
    \item[\rlap{\raisebox{0.3ex}{\hspace{0.4ex}\tiny \ding{52}}}$\square$] für $\omega(t)
    := \gamma(-t)$ gilt $\int_\omega f(s)ds = -\int_\gamma f(s)ds$
    \item[$\square$] falls $f = \nabla g$, dann gilt $\int_\gamma f(s)ds = 0$
    \item[\rlap{\raisebox{0.3ex}{\hspace{0.4ex}\tiny \ding{52}}}$\square$] sei $a = 0, b = 1$
    und $\omega(t) = \gamma(2t), t \in [0, 1/2); \gamma(1), t \in [1/2, 1]$ dann gilt
    $\int_\omega f(s)ds = \int_\gamma f(s)ds$
\end{itemize}

\hrulefill

Sei $X \subseteq \R^n$ offen, $f : X \mapsto \R^n$ ein $C^1$ Funktion, $\gamma: [a,b] \mapsto 
X$ ein parametrisierter Weg. Welche Aussagen sind korrekt?
\begin{itemize}
    \item[\rlap{\raisebox{0.3ex}{\hspace{0.4ex}\tiny \ding{52}}}$\square$] falls $g: X 
    \mapsto \R$ ein Potential von $f$ ist, dann ist für jede Konstante $C \in \R$ $h := g + X$
    auch ein Potential von $f$
    \item[\rlap{\raisebox{0.3ex}{\hspace{0.4ex}\tiny \ding{52}}}$\square$] das Vektorfeld $f$ 
    ist genau dann konservativ, wenn $f$ ein Potential $g$ besitzt
    \item[$\square$] falls $X$ sternförmig ist, ist $f$ konservativ
    \item[$\square$] falls für alle $i,j \in \{1, ..., n\}$ die Gleichung $\partial_{x_j}f_i
    = \partial_{x_i}f_j$ gilt, dann ist $f$ konservativ
    \item[$\square$] seien $A_1, ..., A_m \subseteq X$ offen mit $U_{k = 1}^m A_k = X$. Falls
    $f_{|A_k}$ für alle $k$ konservativ ist, dann ist $f$ konservativ
\end{itemize}

\hrulefill

Seien $A, B \subseteq \R^2$ zwei Mengen und $\X_A, \X_B: \R^2 \mapsto \R$ die dazugehörigen
Charakteristischen Funktionen. Welche Aussagen sind korrekt?
\begin{itemize}
    \item[\rlap{\raisebox{0.3ex}{\hspace{0.4ex}\tiny \ding{52}}}$\square$] $\X_{A \cap B} =
    \X_A \cdot \X_B$
    \item[$\square$] $\X_{A \cup B} = \X_A + \X_B$
    \item[\rlap{\raisebox{0.3ex}{\hspace{0.4ex}\tiny \ding{52}}}$\square$] $\X_{A 
    \backslash B} = \X_A - \X_B$ falls $B \subseteq A$
    \item[$\square$] $\X_A \leq \X_B$ falls $B \subseteq A$
    \item[\rlap{\raisebox{0.3ex}{\hspace{0.4ex}\tiny \ding{52}}}$\square$] $X_{A \cup B} =
    \X_A + \X_B - \X_A \cdot \X_B$
\end{itemize}

\hrulefill

Sei $\R = [a,b] \cross [c,d] \subseteq \R^2, B \subseteq A \subseteq R$ und $f: R 
\mapsto \R$ beschränkt. Welche Aussagen sind korrekt?
\begin{itemize}
    \item[\rlap{\raisebox{0.3ex}{\hspace{0.4ex}\tiny \ding{52}}}$\square$] die Funktion 
    $g(x,y) = e^x$ auf $R$ ist integrierbar
    \item[\rlap{\raisebox{0.3ex}{\hspace{0.4ex}\tiny \ding{52}}}$\square$] falls $f$
    positiv und integrierbar auf $A$ ist, dann gilt $\int_A f(x,y) d(x,y) \geq 0$
    \item[$\square$] falls $f$ auf $B$ und $A$ integrierbar ist, dann $\int_B
    f(x,y)d(x,y) \leq \int_A f(x,y)d(x,y)$
    \item[\rlap{\raisebox{0.3ex}{\hspace{0.4ex}\tiny \ding{52}}}$\square$] es gilt
    $\int_R \X_{A \backslash B}(x,y)d(x,y) = \int_A d(x,y) - \int_B d(x,y)$
\end{itemize}

\hrulefill

Sei $\varphi: D \mapsto \R^2$. Welche Aussagen sind korrekt?
\begin{itemize}
    \item[\rlap{\raisebox{0.3ex}{\hspace{0.4ex}\tiny \ding{52}}}$\square$] für 
    $D = [-2, 2]$ und $\varphi$ Lipschitz, ist Bild $\varphi{[0,1]}$ eine Nullmenge
    \item[$\square$] die Menge $\{(x, e^x) | x \in [0,1]\} \subseteq [0,1]^2$ 
    ist keine Nullmenge
    \item[\rlap{\raisebox{0.3ex}{\hspace{0.4ex}\tiny \ding{52}}}$\square$] Seien
    $\phi: \R^2 \mapsto \R = D$ und $\varphi$ stetig, dann ist die 
    Verkettung $\varphi \circ \phi$ integrierbar auf $[-1, 1]^2$
\end{itemize}

\hrulefill

Welche Aussagen sind korrekt?
\begin{itemize}
    \item[$\square$] $\int_0^\infty e^{-x^2}dx = \pi / 2$
    \item[$\square$] der Satz von Green besagt $\int \int_A F(s)ds = \int_{\partial A}
    (\partial_x f_2 - \partial_y f_1)d(x,y)$
    \item[$\square$] die Jacobi Matrix der Abbildung $\Phi(a,b) = (a^2, ab)^\top$ ist für alle
    $a,b \in \R$ invertierbar
    \item[\rlap{\raisebox{0.3ex}{\hspace{0.4ex}\tiny \ding{52}}}$\square$] die Vektoren $v_1 
    = (0, 1)^\top$ und $v_2 = (1, 0)^\top$ definieren keine positiv orientierte Basis von
    $\R^2$
\end{itemize}

\hrulefill

Sei $f: \R^2 \mapsto \R$. Die Aussagen $\lim_{(x,y) \mapsto (0,0)} f(x,y)$ ist
gleich bedeutend mit:
\begin{itemize}
    \item[\rlap{\raisebox{0.3ex}{\hspace{0.4ex}\tiny \ding{52}}}$\square$]
    $\forall \delta > 0, \exists \epsilon > 0$, so dass $||(x,y)||< \epsilon 
    \Rightarrow |f(x,y)| < \delta$
    \item[$\square$] $\forall \delta > 0, \exists \epsilon > 0$, so dass 
    $|f(x,y)|< \epsilon \Rightarrow ||(x,y)|| < \delta$
    \item[$\square$] $\lim_{x \to 0} f(x,0) = \lim_{y \to 0}f(0,y)$
\end{itemize}

\hrulefill

Es existiert eine reellwertige Funktion $f$ $C^2$, so dass 
$f'' + 27 f' - \pi f = e^{x^2 - x}$ für alle $x \in \R$. \textbf{Wahr}

\hrulefill

Es gibt eine eindeutige Lösung $f$ für das ODE: $y'' + (x^2 + 1)y' + y = 0$, 
$(-1) = -1$. \textbf{Falsch}

\hrulefill

Für eine stetige Funktion $f$ gilt: 
$\int_{[0, 2] \cross [0, 3]} f(x,y)d(x,y) = 6 \int_0^1 \int_0^1 f(2x, 3y)dxdy$.
\textbf{Wahr}

\hrulefill

Wenn $f_1, f_2$ Lösungen der ODE $y'' - xy' + y = \cos(x)$ sind, dann ist auch 
$f_1 + 2 \cdot f_2$ eine Lösung. \textbf{Falsch}

\hrulefill

Wenn $f$ $C^2$ maximal am Punkt $(x_0, y_0)$ ist, dann gilt $\partial^2_{x^2}f = 0.$
\textbf{Falsch}

\hrulefill

Sei $f$ ein Vektorfeld auf $\R^2 - \{0\}$ und $\int_\gamma f(s)ds = 0$ für alle 
geschlossenen Kurven $\gamma$, dann ist $f$ konservativ. \textbf{Wahr}

\hrulefill

Die Funktion $f : \R^2 \mapsto \R, (x,y) \mapsto |xy|$ ist differenzierbar im Punkt (0,0).
\textbf{Wahr}

